\documentclass[11pt]{article}
\usepackage[a4paper,margin=1.9cm]{geometry}
\usepackage[T1]{fontenc}
\usepackage{textcomp}
\usepackage[spanish]{babel}
\usepackage{amsmath,amssymb}
\usepackage{lmodern}
\renewcommand{\familydefault}{\sfdefault}
\usepackage{enumitem}

\newcommand{\vect}[2]{\begin{pmatrix}#1\\#2\end{pmatrix}}
\newcommand{\vectiii}[3]{\begin{pmatrix}#1\\#2\\#3\end{pmatrix}}
\usepackage{tikz}
\usetikzlibrary{calc,arrows.meta,patterns,decorations.pathmorphing}
\usepackage[table]{xcolor}
\usepackage{multicol}
\usepackage{parskip}

\definecolor{unalblue}{RGB}{31,73,124}

\newlist{opciones}{enumerate}{1}
\setlist[opciones]{label=(\alph*),itemsep=6pt,topsep=6pt,leftmargin=1.8em,font=\normalfont}
\setlist[enumerate,1]{itemsep=1.4em,topsep=1em}

\newcommand{\examheader}{%
  \noindent
  \begin{minipage}[t]{0.6\linewidth}
    {\small\bfseries Geometría Vectorial y Analítica --- Parcial 3}\\
    {\small Semestre 2014--01}
  \end{minipage}%
  \hfill
  \begin{minipage}[t]{0.35\linewidth}
    \raggedleft\small Escuela de Matemáticas. Universidad Nacional de Colombia.
  \end{minipage}\\[2pt]
  \rule{\linewidth}{0.6pt}
}
\newcommand{\examfooter}{%
  \vfill
  \rule{\linewidth}{0.4pt}\\[1pt]
  {\scriptsize\textit{Transcripción digital --- MathUNAL}\hfill\textcolor{unalblue}{\textbf{\thepage}}}
}
\pagestyle{empty}

\newcommand{\nota}[1]{{\small\itshape\color{unalblue!70!black} #1}}

\begin{document}
\examheader

\begin{center}
Universidad Nacional de Colombia\\
Facultad de Ciencias --- Escuela de Matemáticas\\
\textbf{Tercer examen parcial de Geometría Vectorial}\\
9 de junio de 2014
\end{center}

\vspace{10pt}
\textbf{PUNTAJE. SOLO PARA USO OFICIAL}

\begin{center}
\begin{tabular}{|c|c|c|c|c|c|}
\hline
1 & 2 & 3 & 4 & 5 & TOTAL \\
\hline
 & & & & & \\
\hline
\end{tabular}
\end{center}

\vspace{10pt}
\textbf{Instrucciones:} No está permitido el uso de calculadora ni de otros aparatos electrónicos. Solucione las preguntas en los espacios en blanco asignados a cada una de ellas. Toda respuesta tiene que estar respaldada por un procedimiento para ser valorada. La última página (página 6) sólo puede usarse para operaciones en borrador; nada de lo escrito en esta página se tendrá en cuenta como parte de la solución del examen. \textbf{Duración del examen: 1 hora y 50 minutos.} No está permitido sacar hojas en blanco ni ningún tipo de apuntes durante el examen.

\begin{enumerate}
\item (20 puntos) Halle una ecuación vectorial para la recta de intersección de los planos $\mathcal{P}_1: x+4y-z=2$ y $\mathcal{P}_2: 3x-y+4z=7$.

\item (20 puntos) Encuentre una ecuación para la elipse que tiene un vértice en el punto $V=\vect{2}{-1}$, centro en el punto $C=\vect{2}{3}$ y excentricidad $e=0.75$.

\item Considere las rectas $\mathcal{L}_1$ y $\mathcal{L}_2$ cuyas ecuaciones paramétricas son:
$$\mathcal{L}_1:\begin{cases}x=-1+2t,\\ y=3,\\ z=-t,\end{cases} t\in\mathbb{R}\qquad\qquad \mathcal{L}_2:\begin{cases}x=2-2r,\\ y=-2,\\ z=r,\end{cases} r\in\mathbb{R}$$

\begin{enumerate}[label=(\alph*)]
\item (10 puntos) Pruebe que $\mathcal{L}_1$ y $\mathcal{L}_2$ son paralelas.
\item (10 puntos) Encuentre una ecuación general del plano $\mathcal{P}$ que contiene a $\mathcal{L}_1$ y $\mathcal{L}_2$.
\end{enumerate}

\item (20 puntos) Pruebe que los puntos $A=\vectiii{2}{-3}{4}$, $B=\vectiii{0}{1}{2}$, $C=\vectiii{-1}{2}{0}$ y $D=\vectiii{2}{-1}{3}$ no son coplanares (es decir, no están situados en un mismo plano) y calcule el volumen del tetraedro con vértices $A$, $B$, $C$ y $D$.

\item Responda \textbf{Falso} o \textbf{Verdadero} a cada uno de los siguientes enunciados y justifique cada respuesta.

\begin{enumerate}[label=(\alph*)]
\item (5 puntos) \textbf{No existen} valores de $t$ para los cuales los vectores $X=\vectiii{1}{-t}{1}$ y $Y=\vectiii{t}{t}{1}$ son linealmente dependientes.
\item (5 puntos) La proyección de todo punto del eje $z$ sobre el plano con ecuación general $x-z=0$ está sobre la recta generada por el vector $U=\vectiii{1}{0}{1}$.
\item (5 puntos) La reflexión de todo punto del eje $z$ con respecto al plano con ecuación general $x-z=0$ está sobre el eje $x$.
\item (5 puntos) Existe un vector $X$ en el espacio tal que sus cosenos directores son los números $\dfrac{1}{3}$, $\dfrac{1}{2}$ y $-\dfrac{1}{5}$.
\end{enumerate}

\end{enumerate}

\examfooter
\end{document}
